% Compact Supplementary Table S2a--b: complete 57-tool library.
% Generated from Supplementary_Table_S2_Assessment_Tool_Library.tex.
% Required packages: booktabs, longtable, array, ragged2e, pdflscape.

\begin{landscape}
\clearpage
\begingroup
\renewcommand{\tablename}{Supplementary Table}
\renewcommand{\thetable}{S2a}
\setlength{\tabcolsep}{3pt}
\setlength{\LTpre}{0pt}
\setlength{\LTpost}{0pt}
\renewcommand{\arraystretch}{1.06}
\scriptsize
\begin{longtable}{@{}
>{\centering\arraybackslash}p{0.50cm}
>{\RaggedRight\arraybackslash}p{4.40cm}
>{\RaggedRight\arraybackslash}p{2.25cm}
>{\RaggedRight\arraybackslash}p{5.90cm}
>{\RaggedRight\arraybackslash}p{4.60cm}
>{\RaggedRight\arraybackslash}p{3.55cm}
@{}}
\caption{Canonical post-stroke assessment-tool library: clinical characteristics, scoring outputs, and evidence sources.}
\label{tab:supplementary_tool_library_clinical}\\
\toprule
\textbf{ID} & \textbf{Tool} & \textbf{COA / mode} & \textbf{Clinical construct and application} & \textbf{Scoring / output} & \textbf{Evidence / source} \\
\midrule
\endfirsthead
\multicolumn{6}{l}{\textit{Supplementary Table S2a continued}}\\
\toprule
\textbf{ID} & \textbf{Tool} & \textbf{COA / mode} & \textbf{Clinical construct and application} & \textbf{Scoring / output} & \textbf{Evidence / source} \\
\midrule
\endhead
\midrule
\multicolumn{6}{r}{\textit{Continued on the next page}}\\
\endfoot
\bottomrule
\endlastfoot
1 & 10-Metre Walk Test (10MWT) & PerfO; performance & Walking speed; Comfortable/fast gait speed over 10 m & Gait speed in m/s from a timed 10-m walk & G \cite{S1VADoD2024StrokeRehab,S1CSBPR2025RehabPart2} \\
2 & 2-Minute Walk Test (2MWT) & PerfO; performance & Walking capacity; Short-distance walking endurance/capacity & Distance walked in metres over 2 min & S \cite{S22MinuteWalkTest22502805} \\
3 & 6-Minute Walk Test (6MWT) & PerfO; performance & Walking endurance; Distance walked in 6 min & Distance walked in metres over 6 min & G \cite{S1VADoD2024StrokeRehab,S1CSBPR2025RehabPart2} \\
4 & Action Research Arm Test (ARAT) & ClinRO; task-based rating & Upper-limb activity/function; Grasp, grip, pinch, gross arm movement & 19 tasks scored 0--3; total 0--57 & G \cite{S1Winstein2016AHAASARehab,S1VADoD2024StrokeRehab} \\
5 & Berg Balance Scale (BBS) & ClinRO; task-based rating & Standing balance; Static and dynamic standing balance/fall risk & 14 tasks scored 0--4; total 0--56 & G \cite{S1Winstein2016AHAASARehab,S1VADoD2024StrokeRehab} \\
6 & Box and Block Test (BBT) & PerfO; performance & Manual dexterity (gross); Number of blocks moved across a partition in 60 s & Number of blocks transferred in 60 s & G \cite{S1Winstein2016AHAASARehab,S1VADoD2024StrokeRehab} \\
7 & Brunel Balance Assessment (BBA) & PerfO; performance & Hierarchical balance; Hierarchical balance assessment across sitting, standing, and stepping & 12 hierarchical balance tasks; highest level achieved & S \cite{S2BrunelBalanceAsses15573837} \\
8 & Chedoke Arm and Hand Activity Inventory (CAHAI) & ClinRO; task-based rating & Bilateral functional arm/hand activity; Performance of bimanual everyday tasks & Functional tasks scored 1--7; version-dependent total & G \cite{S1Winstein2016AHAASARehab,S1CSBPR2025RehabPart2} \\
9 & Community Balance \& Mobility Scale (CB\&M) & ClinRO; task-based rating & Community balance and mobility; High-level balance and community mobility in higher-functioning ambulators & 13 performance items; total 0--96 & S \cite{S2CommunityBalanceMo20510980} \\
10 & Dynamic Gait Index (DGI) & ClinRO; task-based rating & Complex/dynamic walking; Gait under changing task demands & Eight tasks scored 0--3; total 0--24 & G \cite{S1CSBPR2025RehabPart2} \\
11 & Elderly Mobility Scale (EMS) & ClinRO; task-based rating & Functional mobility; Functional mobility including locomotion, balance, and transfers & Seven items; total 0--20 & S \cite{S2ElderlyMobilitySca9408676} \\
12 & Five Times Sit-to-Stand (5TSTS) & PerfO; performance & Lower-limb strength/transfer; Speed of 5 sit-to-stand repetitions & Time in seconds for five repetitions & G \cite{S1StrokeFoundationLivingGuidelineRehab} \\
13 & Frenchay Arm Test (FAT) & PerfO; performance & Upper-limb activity; Gross arm and hand functional tasks & Five dichotomous tasks; total 0--5 & S \cite{S2FrenchayArmTest3612152} \\
14 & Function In Sitting Test (FIST) & ClinRO; task-based rating & Sitting balance; Functional sitting balance and seated postural control & 14 tasks scored 0--4; total 0--56 & S \cite{S2FunctionInSittingT20716989} \\
15 & Functional Gait Assessment (FGA) & ClinRO; task-based rating & Complex/dynamic walking; Gait under changing demands (head turns, pivot, obstacle, narrow base, eyes closed) & 10 tasks scored 0--3; total 0--30 & G \cite{S1AltMurphy2025ESOMotorRehab,S1Johnston2021AFOFESPostStroke} \\
16 & Functional Reach Test (FRT) & PerfO; performance & Anticipatory postural control; Max forward reach in standing & Mean forward reach of three trials, in cm or inches & G \cite{S1Winstein2016AHAASARehab,S1StrokeFoundationLivingGuidelineRehab} \\
17 & Jebsen-Taylor Hand Function Test (JTHFT) & PerfO; performance & Hand function and dexterity; Timed hand function across daily manual tasks & Seven timed hand-function subtests; output in seconds & S \cite{S2JebsenTaylorHandFu15053129} \\
18 & Manual Function Test (MFT) & PerfO; performance & Upper-limb motor function; Upper-limb motor/manual function including arm and hand items & Eight tasks; motor-function score 0--32 & S \cite{S2ManualFunctionTest19106794} \\
19 & Mini-BESTest & ClinRO; task-based rating & Dynamic balance (multi-system); Anticipatory, reactive, sensory-orientation and dynamic-gait balance & 14 items scored 0--2; total 0--28 & G \cite{S1CSBPR2025RehabPart2} \\
20 & Modified Emory Functional Ambulation Profile (mEFAP) & PerfO; summed task times & Timed functional ambulation over varied terrain; Timed walking across 5 environmental tasks (floor, carpet, up-and-go, obstacles, stairs) & Five timed terrain tasks; times summed without weighting, manual assistance recorded separately; lower is better & G \cite{S1Johnston2021AFOFESPostStroke} \\
21 & Nine-Hole Peg Test (9HPT) & PerfO; performance & Manual dexterity (fine); Time to place and remove 9 pegs & Completion time in seconds per hand & G \cite{S1VADoD2024StrokeRehab,S1CSBPR2025RehabPart2} \\
22 & Short Physical Performance Battery (SPPB) & PerfO; performance & Lower-limb physical performance; Lower-extremity physical performance: gait speed, balance, chair stand & Three components scored 0--4; total 0--12 & S \cite{S2ShortPhysicalPerfo40162876} \\
23 & Six-Spot Step Test (SSST) & PerfO; performance & Complex walking and dynamic balance; Dynamic balance and agility during walking with step targets & Course-completion time in seconds & S \cite{S2SixSpotStepTest34128605} \\
24 & Stroke Upper Limb Capacity Scale (SULCS) & PerfO; performance & Upper-limb capacity; Hierarchical stroke upper-limb capacity from basic to complex tasks & Ten dichotomous tasks; total 0--10 & S \cite{S2StrokeUpperLimbCap21272717} \\
25 & Timed Up and Go (TUG) & PerfO; performance & Functional mobility/dynamic balance; Time to rise, walk 3 m, turn, return, sit & Total completion time in seconds & G \cite{S1Winstein2016AHAASARehab,S1VADoD2024StrokeRehab} \\
26 & Tinetti/POMA & ClinRO; task-based rating & Balance+gait/fall risk; Observed balance+gait subscales & 16 balance and gait items; total 0--28 & G \cite{S1Winstein2016AHAASARehab} \\
27 & Trunk Control Test (TCT) & ClinRO; task-based rating & Trunk control; 4 trunk movements & Four weighted items; total 0--100 & G \cite{S1CSBPR2025RehabPart2,S1StrokeFoundationLivingGuidelineRehab} \\
28 & Trunk Impairment Scale (TIS, Verheyden) & ClinRO; task-based rating & Sitting balance/trunk control; Static and dynamic sitting balance+trunk coordination & 17 items; total 0--23 & G \cite{S1CSBPR2025RehabPart2,S1StrokeFoundationLivingGuidelineRehab} \\
29 & Wolf Motor Function Test (WMFT) & PerfO/ClinRO; timed and strength outputs PerfO, FAS rating ClinRO & Upper-limb motor function and speed; Timed+rated functional arm tasks & 15 timed tasks plus functional-ability ratings 0--5; grip strength (task 14) and weight to box (task 7) recorded separately & G \cite{S1Winstein2016AHAASARehab,S1VADoD2024StrokeRehab} \\
30 & Ability for Basic Movement Scale II (ABMS II) & ClinRO; task-based rating & Basic mobility; Assistance needed for five basic movements from supine lying to standing & Five items graded 1--6; total 5--30 & S \cite{S2AbilityforBasicMov20140415} \\
31 & Barthel Index (BI) & ClinRO; observation/interview & Basic ADL (incl. mobility, transfers, stairs); 10 ADL items including transfers, ambulation, stairs & Ten-item weighted total; 0--20 or 0--100 when rescaled & G \cite{S1Winstein2016AHAASARehab,S1VADoD2024StrokeRehab} \\
32 & Brunnstrom Recovery Stages (BRS) & ClinRO; clinician & Motor recovery stage; Stage-based motor recovery and synergy pattern after stroke & Six ordinal recovery stages (I-VI) & S \cite{S2BrunnstromRecovery27495103} \\
33 & Chedoke-McMaster Stroke Assessment (CMSA) & ClinRO; clinician & Impairment+activity (incl. leg, foot, postural control, walking); Stages of motor recovery+activity (incl. walking) & Impairment stages 1--7 plus a 100-point Activity Inventory & G \cite{S1Winstein2016AHAASARehab,S1CSBPR2025RehabPart2} \\
34 & Composite Spasticity Index (CSI) & ClinRO; clinician & Spasticity; Composite lower-limb spasticity outcome & Tendon jerk, passive resistance and clonus; total 0--16 & S \cite{S2CompositeSpasticit38907596} \\
35 & Fugl-Meyer LE (FMA-LE) & ClinRO; clinician & Sensorimotor impairment (LE motor and balance subscales gait-relevant); Motor recovery (UE+LE), balance, sensation, ROM & 0--2 per item; motor subscore 0--34 & G \cite{S1Winstein2016AHAASARehab,S1VADoD2024StrokeRehab} \\
36 & Fugl-Meyer UE (FMA-UE) & ClinRO; clinician & Upper-limb motor impairment; Reflexes, volitional movement, coordination/speed of the paretic arm & 0--2 per item; motor subscore 0--66 & G \cite{S1VADoD2024StrokeRehab,S1CSBPR2025RehabPart2} \\
37 & Functional Ambulation Category (FAC) & ClinRO; task-based rating & Walking independence; Level of assistance needed to walk & Ordinal category 0--5 based on walking assistance & G \cite{S1Winstein2016AHAASARehab,S1VADoD2024StrokeRehab} \\
38 & Functional Independence Measure (FIM) & ClinRO; observation/interview & ADL independence (motor including locomotion/transfer); 18 items (13 motor including locomotion+5 cognitive) & 18 items scored 1--7; total 18--126 & G \cite{S1Winstein2016AHAASARehab,S1VADoD2024StrokeRehab} \\
39 & Medical Research Council (MRC) muscle scale & ClinRO; clinician & Muscle strength; Manual muscle strength grading & Ordinal strength grade 0--5 per muscle & S \cite{S2MedicalResearchCou40493483} \\
40 & Modified Ashworth Scale (MAS) & ClinRO; clinician & Muscle spasticity (affects gait); Resistance to passive movement (tone) & Ordinal grades 0, 1, 1+, 2, 3 and 4 & G \cite{S1Winstein2016AHAASARehab,S1VADoD2024StrokeRehab} \\
41 & modified Rankin Scale (mRS) & ClinRO; clinician & Global disability/dependence; Overall disability; grades 2--4 defined by walking/assistance & Single global disability grade 0--6 & G \cite{S1Winstein2016AHAASARehab,S1VADoD2024StrokeRehab} \\
42 & Modified Tardieu Scale (MTS) & ClinRO; clinician & Muscle spasticity (velocity-dependent); Spasticity vs contracture at different velocities & Quality of muscle reaction plus R1/R2 joint angles & S \cite{S2RehabMeasuresModifiedTardieu} \\
43 & Motor Assessment Scale (MoAS) & ClinRO; task-based rating & Motor function (incl. walking item); 8 functional motor items including walking, sitting balance & Eight scored items, 0--6 each; total 0--48; eight-item version; general tonus of the original version not mapped & G \cite{S1StrokeFoundationLivingGuidelineRehab} \\
44 & Motricity Index (MI) & ClinRO; clinician & Voluntary muscle strength (incl. leg); Strength of key arm and leg movements & Weighted ordinal strength scores; limb index 0--100 & G \cite{S1Winstein2016AHAASARehab,S1CSBPR2025RehabPart2} \\
45 & NIH Stroke Scale (NIHSS) & ClinRO; clinician & Stroke severity (incl. limb-motor and limb-ataxia items); Neuro deficit severity; arm/leg motor+limb-ataxia items & Item-level ordinal ratings; total 0--42 & G \cite{S1Winstein2016AHAASARehab,S1StrokeFoundationLivingGuidelineRehab} \\
46 & Scale for Contraversive Pushing (SCP) & ClinRO; clinician & Postural orientation and pushing; Pusher behavior and postural orientation after stroke & Three domains assessed in sitting and standing; total 0--6 & S \cite{S2ScaleforContravers17082502} \\
47 & Shoulder Abduction Finger Extension (SAFE) score & ClinRO; clinician (MRC) & Upper-limb motor prognosis; Early upper-limb motor predictor based on shoulder abduction and finger extension strength & Two MRC strength grades; total 0--10 & S \cite{S2ShoulderAbductionF20167916} \\
48 & STREAM & ClinRO; task-based rating & Voluntary movement+basic mobility; Limb movement+mobility (incl. walking) items & 30 items scored 0--2 or 0--3; total 0--70 & G \cite{S1Winstein2016AHAASARehab,S1Johnston2021AFOFESPostStroke} \\
49 & Stroke Impairment Assessment Set (SIAS) & ClinRO; clinician & Multidomain stroke impairment; Multi-domain stroke impairment including motor, trunk, sensory, and ROM domains & 22 items with item-specific ordinal ratings; total 0--76 & S \cite{S2StrokeImpairmentAs10895996} \\
50 & ABILHAND & PRO; self-report & Perceived manual ability; Patient-perceived manual ability in bimanual daily tasks & 23 activities rated impossible, difficult or easy; Rasch-derived ability estimate & S \cite{S2ABILHAND11441211} \\
51 & Activities-specific Balance Confidence (ABC) & PRO; self-report & Balance self-efficacy; Confidence in balance across 16 activities & 16 confidence ratings, 0--100\%; mean score & G \cite{S1Winstein2016AHAASARehab,S1VADoD2024StrokeRehab} \\
52 & Activity Measure for Post-Acute Care (AM-PAC) & ClinRO; clinician-completed 6-Clicks short form & Post-acute mobility and function; Basic Mobility inpatient short form (6-Clicks): six activities rated for difficulty or help needed & Six items rated 1--4 (raw 6--24), converted to a standardized scale score & S \cite{S2ActivityMeasurefor22343646} \\
53 & Falls Efficacy Scale (FES) & PRO; self-report & Falls self-efficacy; Fear of falling and fall-related self-efficacy & Ten self-report items; summed score 10--100 & S \cite{S2FallsEfficacyScale10588538} \\
54 & Falls Efficacy Scale-International (FES-I) & PRO; self-report & Concern about falling; Concern about falling during basic and more demanding physical and social activities & Sixteen items rated 1--4; total 16--64 & S \cite{S2FallsEfficacyScale33992239} \\
55 & Motor Activity Log (MAL) & PRO; self-report & Real-world arm use (amount and quality); Self-report of how much/how well the paretic arm is used in daily tasks & 30-item version: amount-of-use and quality-of-movement ratings per activity, each 0--5; scale scores are item means & G \cite{S1Winstein2016AHAASARehab,S1VADoD2024StrokeRehab} \\
56 & Stroke Impact Scale (SIS) & PRO; self-report & Self-reported impact including mobility domain; 8 domains including mobility and hand/strength & 59 items; domain scores transformed to 0--100 & G \cite{S1Winstein2016AHAASARehab,S1VADoD2024StrokeRehab} \\
57 & Rivermead Mobility Index (RMI) & PRO/PerfO; 14 self-reported items, 1 observed item & Functional mobility; Bed mobility, transfers, walking, stairs & 15 dichotomous items; total 0--15 & G \cite{S1CSBPR2025RehabPart2,S1AltMurphy2025ESOMotorRehab} \\
\end{longtable}
\noindent\scriptsize\textit{Notes.} G, guideline evidence stream; S, supplementary clinical-literature stream. COA, clinical outcome assessment; ClinRO, clinician-reported outcome; PerfO, performance outcome; PRO, patient-reported outcome. Scoring ranges refer to the listed canonical form and may differ across modified versions. Evidence/source identifies the local evidence stream and key report used for extraction.
\endgroup

\end{landscape}

\clearpage
\begingroup
\renewcommand{\tablename}{Supplementary Table}
\renewcommand{\thetable}{S2b}
\setlength{\tabcolsep}{2.2pt}
\setlength{\LTpre}{0pt}
\setlength{\LTpost}{0pt}
\renewcommand{\arraystretch}{1.06}
\newcommand{\mcmark}{\(\bullet\)}
\small
\begin{longtable}{@{}
>{\centering\arraybackslash}p{0.6cm}
>{\RaggedRight\arraybackslash}p{2.3cm}
>{\centering\arraybackslash}p{0.8cm}
*{7}{>{\centering\arraybackslash}p{0.75cm}}
>{\centering\arraybackslash}p{1.9cm}
@{}}
\caption{Movement-component coverage and video availability of score-defining information for the 57 canonical post-stroke assessment tools.}
\label{tab:supplementary_tool_library_motion}\\
\toprule
\multirow{2}{*}{\textbf{ID}} &
\multirow{2}{*}{\textbf{Scale}} &
\multirow{2}{*}{\textbf{Unit}} &
\multicolumn{7}{c}{\textbf{Movement components}} &
\multirow{2}{*}{\shortstack{\textbf{Video}\\\textbf{F/P/N}}} \\
\cmidrule(lr){4-10}
& & & \textbf{ROM} & \textbf{S/F} & \textbf{Coord.} & \textbf{Bal.} & \textbf{Dex.} & \textbf{Gait} & \textbf{Q/C} & \\
\midrule
\endfirsthead
\multicolumn{11}{l}{\textit{Supplementary Table S2b continued}}\\
\toprule
\textbf{ID} & \textbf{Scale} & \textbf{Unit} & \textbf{ROM} & \textbf{S/F} & \textbf{Coord.} & \textbf{Bal.} & \textbf{Dex.} & \textbf{Gait} & \textbf{Q/C} & \textbf{Video} \\
\midrule
\endhead
\midrule
\multicolumn{11}{r}{\textit{Continued on the next page}}\\
\endfoot
\bottomrule
\endlastfoot
1 & 10MWT & T &  &  &  &  &  & \mcmark &  & F1/P0/N0 \\
2 & 2MWT & T &  &  &  &  &  & \mcmark &  & F1/P0/N0 \\
3 & 6MWT & T &  &  &  &  &  & \mcmark &  & F1/P0/N0 \\
4 & ARAT & S &  &  & \mcmark &  & \mcmark &  & \mcmark & F4/P0/N0 \\
5 & BBS & I &  &  &  & \mcmark &  & \mcmark & \mcmark & F10/P4/N0 \\
6 & BBT & T &  &  &  &  & \mcmark &  &  & F1/P0/N0 \\
7 & BBA & S &  &  & \mcmark & \mcmark &  & \mcmark & \mcmark & F3/P0/N0 \\
8 & CAHAI & T &  &  & \mcmark &  & \mcmark &  &  & F0/P1/N0 \\
9 & CB\&M & T &  &  & \mcmark & \mcmark &  & \mcmark & \mcmark & F1/P0/N0 \\
10 & DGI & I &  &  & \mcmark & \mcmark &  & \mcmark & \mcmark & F6/P2/N0 \\
11 & EMS & I &  &  &  & \mcmark &  & \mcmark &  & F4/P3/N0 \\
12 & 5TSTS & T &  &  &  &  &  & \mcmark &  & F1/P0/N0 \\
13 & FAT & T & \mcmark &  & \mcmark &  & \mcmark &  &  & F1/P0/N0 \\
14 & FIST & T & \mcmark &  & \mcmark & \mcmark &  &  & \mcmark & F0/P1/N0 \\
15 & FGA & I &  &  & \mcmark & \mcmark &  & \mcmark & \mcmark & F9/P1/N0 \\
16 & FRT & T &  &  &  & \mcmark &  &  &  & F1/P0/N0 \\
17 & JTHFT & T &  &  &  &  & \mcmark &  &  & F1/P0/N0 \\
18 & MFT & I & \mcmark & \mcmark & \mcmark &  & \mcmark &  &  & F8/P0/N0 \\
19 & Mini-BESTest & I & \mcmark &  &  & \mcmark &  & \mcmark & \mcmark & F13/P1/N0 \\
20 & mEFAP & T &  &  &  &  &  & \mcmark &  & F1/P0/N0 \\
21 & 9HPT & T &  &  &  &  & \mcmark &  &  & F1/P0/N0 \\
22 & SPPB & S &  &  &  & \mcmark &  & \mcmark &  & F3/P0/N0 \\
23 & SSST & T &  &  &  &  &  & \mcmark &  & F1/P0/N0 \\
24 & SULCS & T & \mcmark &  & \mcmark &  & \mcmark &  &  & F1/P0/N0 \\
25 & TUG & T &  &  &  &  &  & \mcmark &  & F1/P0/N0 \\
26 & POMA & I &  &  & \mcmark & \mcmark &  & \mcmark & \mcmark & F16/P0/N0 \\
27 & TCT & I & \mcmark &  & \mcmark & \mcmark &  &  & \mcmark & F4/P0/N0 \\
28 & TIS & I &  &  & \mcmark & \mcmark &  &  & \mcmark & F17/P0/N0 \\
29 & WMFT & S & \mcmark & \mcmark & \mcmark &  & \mcmark &  & \mcmark & F3/P0/N1 \\
30 & ABMS II & I &  &  &  & \mcmark &  & \mcmark & \mcmark & F0/P5/N0 \\
31 & BI\textsuperscript{a} & I &  &  &  &  &  & \mcmark &  & F0/P3/N0 \\
32 & BRS & T & \mcmark & \mcmark & \mcmark &  &  &  & \mcmark & F0/P1/N0 \\
33 & CMSA & T & \mcmark & \mcmark & \mcmark & \mcmark &  & \mcmark & \mcmark & F0/P1/N0 \\
34 & CSI & T &  &  &  &  &  &  &  & F0/P0/N1 \\
35 & FMA-LE\textsuperscript{a} & S & \mcmark & \mcmark & \mcmark &  &  &  & \mcmark & F1/P1/N1 \\
36 & FMA-UE\textsuperscript{a} & S & \mcmark & \mcmark & \mcmark &  & \mcmark &  & \mcmark & F2/P2/N1 \\
37 & FAC & T &  &  &  & \mcmark &  & \mcmark &  & F0/P1/N0 \\
38 & FIM\textsuperscript{a} & T &  &  &  &  &  & \mcmark &  & F0/P1/N0 \\
39 & MRC & T & \mcmark & \mcmark &  &  &  &  &  & F0/P1/N0 \\
40 & MAS & T &  &  &  &  &  &  &  & F0/P0/N1 \\
41 & mRS & T &  &  &  &  &  & \mcmark &  & F0/P1/N0 \\
42 & MTS & T & \mcmark &  &  &  &  &  &  & F0/P1/N0 \\
43 & MoAS & I & \mcmark & \mcmark & \mcmark & \mcmark & \mcmark & \mcmark & \mcmark & F8/P0/N0 \\
44 & MI & T & \mcmark & \mcmark &  &  &  &  &  & F0/P1/N0 \\
45 & NIHSS\textsuperscript{a} & S &  & \mcmark & \mcmark &  &  &  &  & F2/P0/N0 \\
46 & SCP & S &  &  &  & \mcmark &  &  & \mcmark & F2/P0/N1 \\
47 & SAFE & T & \mcmark & \mcmark &  &  &  &  &  & F0/P1/N0 \\
48 & STREAM & S & \mcmark & \mcmark & \mcmark & \mcmark &  & \mcmark & \mcmark & F3/P0/N0 \\
49 & SIAS\textsuperscript{a} & S & \mcmark & \mcmark & \mcmark & \mcmark &  &  & \mcmark & F1/P2/N0 \\
50 & ABILHAND & T &  &  &  &  & \mcmark &  &  & F0/P0/N1 \\
51 & ABC & T &  &  &  & \mcmark &  & \mcmark &  & F0/P0/N1 \\
52 & AM-PAC\textsuperscript{a} & T &  &  &  & \mcmark &  & \mcmark &  & F0/P1/N0 \\
53 & FES & T &  &  &  & \mcmark &  & \mcmark &  & F0/P0/N1 \\
54 & FES-I & T &  &  &  & \mcmark &  & \mcmark &  & F0/P0/N1 \\
55 & MAL & T &  &  &  &  & \mcmark &  & \mcmark & F0/P0/N1 \\
56 & SIS\textsuperscript{a} & T &  & \mcmark &  & \mcmark & \mcmark & \mcmark &  & F0/P0/N1 \\
57 & RMI & I &  &  &  & \mcmark &  & \mcmark &  & F1/P0/N14 \\
\end{longtable}
\noindent\scriptsize\textit{Notes.} Filled circles indicate score-defining movement components at scale level. The first six columns are primary movement components; Q/C (movement quality and compensation) is a cross-component modifier rather than an independent clinical domain. Components were coded with three rules. (i) For items rated against task-specific criteria, including pass/fail criteria, a component was marked when the scored task requires it and the scoring criteria distinguish performance on it. (ii) For tests scored only by time, distance or count, only the component that the measured quantity directly expresses was marked (for example, walking or sit-to-stand time marks gait and functional mobility, and reach distance in standing marks balance); for tests scored by level of assistance, the components named in the assistance criteria were marked. Abilities that a task requires but that do not change the score were not coded. (iii) For patient-reported measures, components referenced by item content were marked. Q/C was marked only when the scoring criteria explicitly grade movement quality, normality of execution, abnormal synergy or the use of a compensatory strategy such as arm support. Passive tone, reflexes and sensation fall outside the six components, and the endurance construct of the 2MWT and 6MWT is represented only by the Gait column. The coding framework was informed by the ICF framework, stroke rehabilitation guidelines and SRRR consensus domains \cite{WHO2001ICF,S1Winstein2016AHAASARehab,S1VADoD2024StrokeRehab,S1NICE2023NG236,CoreKwakkel2017SRRRSensorimotor,CoreKwakkel2019SRRRUpperLimb,CoreVanCriekinge2024SRRRBalanceMobility}. ROM, joint mobility and range of motion; S/F, muscle strength and force production; Coord., coordination and selective motor control; Bal., balance and postural control; Dex., manual dexterity and upper-limb task execution; Gait, gait and functional mobility. Unit gives the mapping unit: I, each item mapped separately; S, subscales, domains or item groups mapped separately; T, the whole tool mapped as one entry. Because units differ, F/P/N counts are not comparable across tools. Video F/P/N gives the number of entries whose score-defining information is fully, partially or not available from ordinary video under an adequate recording: Full, all information used by the scoring rules is visible, including movement, time, counts, distances with a visible reference, and assistance or support that actually occurs; Partial, some grade-defining information is not visible (force or resistance, palpation, the amount of assistance given, verbal cues or responses, judgements without a visible trace, or usual performance outside the recording); None, the score does not come from visible performance (patient report, reflexes, tone or force read by touch or instrument). Counts refer to 196 mapping entries; item-level details and representative machine-observable features are provided in Supplementary Data 1. \textsuperscript{a}Coding restricted to movement-related content: BI, mobility items; FIM, motor items; NIHSS, motor and limb-ataxia items; SIAS, motor, trunk and ROM items; AM-PAC, basic mobility domain; SIS, strength, hand and mobility domains; FMA-UE and FMA-LE, motor items.
\endgroup

